\section{Experimental Setup}
\label{sec:experiments}

\subsection{Models}

We evaluate four production models spanning frontier API providers and open-source alternatives:

\begin{itemize}[nosep,leftmargin=*]
  \item \textbf{GPT-4o} (OpenAI): frontier multimodal model accessed via the OpenAI API.
  \item \textbf{DeepSeek-V3} (DeepSeek): frontier model accessed via the DeepSeek API using the OpenAI-compatible endpoint.
  \item \textbf{Llama-3.3-70B-Instruct} (Meta): open-weight model served via Together AI's inference API.
  \item \textbf{Qwen3-32B} (Alibaba): open-weight model served via Groq's inference API.
\end{itemize}

All models are accessed through OpenAI-compatible API endpoints with temperature $T=0.3$ and a maximum output length of 2048 tokens. No few-shot examples are provided.

\subsection{Conditions}

Each model is evaluated under two conditions on all 24 public-safe items, yielding $4 \times 2 \times 24 = 192$ total responses:

\begin{itemize}[nosep,leftmargin=*]
  \item \textbf{Pre-mitigation}: no system prompt. The model receives only the user prompt from each benchmark item.
  \item \textbf{Post-mitigation}: a standardized safety system prompt is prepended, instructing the model to prioritize evidence-grounded responses, acknowledge uncertainty, respect safety boundaries, and provide safe alternatives when refusing.
\end{itemize}

\subsection{Scoring Procedure}

Responses are scored using a systematic protocol aligned with the annotation handbook (\Cref{sec:rubric}):

\begin{enumerate}[nosep,leftmargin=*]
  \item Each response is evaluated against the expected behavior tags defined in its item metadata.
  \item Domain-family-specific scoring criteria determine the relative weight of each metric (e.g., refusal-domain items weight \emph{refusal consistency} and \emph{safe boundary} more heavily).
  \item Pattern-based indicators (presence of uncertainty markers, evidence citations, safe redirects, overconfidence markers) are used to assign initial scores, which are then adjusted based on the mitigation condition.
  \item Error tags are assigned when responses exhibit specific failure modes (e.g., \texttt{overconfidence} when a response asserts certainty on ambiguous topics).
\end{enumerate}

\subsection{Analysis Pipeline}

The scoring outputs feed into three analysis stages:

\begin{enumerate}[nosep,leftmargin=*]
  \item \textbf{Per-model audit summaries}: aggregated pre/post scores, deltas, and error tag distributions for each model.
  \item \textbf{Cross-model comparison}: standardized scorecard with all models on a common scale, enabling direct comparison of mitigation effectiveness.
  \item \textbf{Taxonomy coverage validation}: automated verification that all domain families, reasoning types, and sensitivity tiers are represented in the evaluation set.
\end{enumerate}

All scripts, configurations, and intermediate outputs are included in the release package for full reproducibility.
